% GENERATED FILE. Edit canonical lessons, then run npm run generate:books.
% canonical-source-hash: fnv1a64:022a89aa466af88b
% canonical-lessons: RU-C485-toshnit, RU-C485-toshnota, RU-C485-golovokruzheniye, RU-C485-obmorok, RU-C485-znobit

\chapter{Feeling Sick and Faint}
\label{ch:ru-feeling-sick-and-faint}
\hlchaptermodality{485}

\begin{chapteropening}
\textbf{By the end of this chapter:} \emph{I can describe feeling sick, dizzy and faint.}

Complete the last lesson of chapter 485: I can describe feeling sick, dizzy and faint.
\end{chapteropening}

\section[\texorpdfstring{\ru{тошнить} (toshnít)}{тошнить (toshnít)}]{\ru{тошнить} (toshnít) — to feel sick (impersonal)}

\label{lesson:RU-C485-toshnit}



\begin{quote}
Before the new one: say the Russian for a strain, then the Russian for to limp.
\end{quote}

\subsection*{You'll want to know: \ru{тошнить}}
\textbf{\ru{тошнить}} (\emph{toshnít}) — \textquotedblleft{}to feel sick (impersonal)\textquotedblright{}.

\begin{quote}
\ru{В} \ru{автобусе} \ru{меня} \ru{часто} \ru{тошнит}.
\end{quote}

\begin{quote}
\emph{V avtóbuse menyá chásto toshnít.}
\end{quote}

\begin{quote}
I often feel sick on the bus.
\end{quote}

Impersonal: \emph{\ru{меня} \ru{тошнит}} — I feel sick. Imperfective; perfective \ru{стошнить} — to be sick.

\begin{cousinweb}
\emph{\ru{тошнить}} is built on \emph{\ru{тошно}} 'sickening, miserable'; EtymDB has no entry, and the root is usually linked to the one in \emph{\ru{тоска}} 'longing'.
\end{cousinweb}

\begin{grammarlens}[title={What you've built}]
One more word for telling, explaining and describing at B1.
\end{grammarlens}

\subsection*{Your turn}
\begin{itemize}
\raggedright
  \item \emph{Say it:} \emph{toshnít}
  \item \emph{Say it:} \emph{toshnít}, once more
  \item \emph{Say it:} a sentence of your own with \emph{toshnít}
  \item \emph{Recall:} say the Russian for a strain, then the Russian for to limp, then say \emph{toshnít} again
\end{itemize}

\begin{tcolorbox}[breakable,colback=teal!4,colframe=teal!35!black,title={Before you move on}]
What does \ru{тошнить} mean? (\textquotedblleft{}To feel sick (impersonal)\textquotedblright{}.) Say it once more.
\end{tcolorbox}
\section[\texorpdfstring{\ru{тошнота} (toshnotá)}{тошнота (toshnotá)}]{\ru{тошнота} (toshnotá) — nausea}

\label{lesson:RU-C485-toshnota}



\begin{quote}
Before the new one: say the Russian for to limp, then the Russian for to feel sick.
\end{quote}

\subsection*{You'll want to know: \ru{тошнота}}
\textbf{\ru{тошнота}} (\emph{toshnotá}) — \textquotedblleft{}nausea\textquotedblright{}.

\begin{quote}
\ru{Это} \ru{лекарство} \ru{помогает} \ru{от} \ru{тошноты}.
\end{quote}

\begin{quote}
\emph{Éto lekárstvo pomogáyet ot toshnotý.}
\end{quote}

\begin{quote}
This medicine helps with nausea.
\end{quote}

Feminine.

\begin{cousinweb}
\emph{\ru{тошнота}} is built on \emph{\ru{тошно}} 'sickening' with \emph{-\ru{ота}}; EtymDB has no entry, and the root is usually linked to \emph{\ru{тоска}}.
\end{cousinweb}

\begin{grammarlens}[title={What you've built}]
One more word for telling, explaining and describing at B1.
\end{grammarlens}

\subsection*{Your turn}
\begin{itemize}
\raggedright
  \item \emph{Say it:} \emph{toshnotá}
  \item \emph{Say it:} \emph{toshnotá}, once more
  \item \emph{Say it:} a sentence of your own with \emph{toshnotá}
  \item \emph{Recall:} say the Russian for to limp, then the Russian for to feel sick, then say \emph{toshnotá} again
\end{itemize}

\begin{tcolorbox}[breakable,colback=teal!4,colframe=teal!35!black,title={Before you move on}]
What does \ru{тошнота} mean? (\textquotedblleft{}Nausea\textquotedblright{}.) Say it once more.
\end{tcolorbox}
\section[\texorpdfstring{\ru{головокружение} (golovokruzhéniye)}{головокружение (golovokruzhéniye)}]{\ru{головокружение} (golovokruzhéniye) — dizziness}

\label{lesson:RU-C485-golovokruzheniye}



\begin{quote}
Before the new one: say the Russian for to feel sick, then the Russian for nausea.
\end{quote}

\subsection*{You'll want to know: \ru{головокружение}}
\textbf{\ru{головокружение}} (\emph{golovokruzhéniye}) — \textquotedblleft{}dizziness\textquotedblright{}.

\begin{quote}
\ru{Летом} \ru{от} \ru{жары} \ru{у} \ru{меня} \ru{было} \ru{головокружение}.
\end{quote}

\begin{quote}
\emph{Létom ot zharý u menyá býlo golovokruzhéniye.}
\end{quote}

\begin{quote}
In the summer the heat made me dizzy.
\end{quote}

Neuter. \emph{\ru{У} \ru{меня} \ru{кружится} \ru{голова}} — I feel dizzy.

\begin{cousinweb}
\emph{\ru{головокружение}} is a compound of \emph{\ru{голова}} 'head', inherited from Proto-Slavic \emph{golva} (EtymDB), and \emph{\ru{кружение}} 'spinning', built on \emph{\ru{круг}} 'circle'.
\end{cousinweb}

\begin{grammarlens}[title={What you've built}]
One more word for telling, explaining and describing at B1.
\end{grammarlens}

\subsection*{Your turn}
\begin{itemize}
\raggedright
  \item \emph{Say it:} \emph{golovokruzhéniye}
  \item \emph{Say it:} \emph{golovokruzhéniye}, once more
  \item \emph{Say it:} a sentence of your own with \emph{golovokruzhéniye}
  \item \emph{Recall:} say the Russian for to feel sick, then the Russian for nausea, then say \emph{golovokruzhéniye} again
\end{itemize}

\begin{tcolorbox}[breakable,colback=teal!4,colframe=teal!35!black,title={Before you move on}]
What does \ru{головокружение} mean? (\textquotedblleft{}Dizziness\textquotedblright{}.) Say it once more.
\end{tcolorbox}
\section[\texorpdfstring{\ru{обморок} (óbmorok)}{обморок (óbmorok)}]{\ru{обморок} (óbmorok) — a faint}

\label{lesson:RU-C485-obmorok}



\begin{quote}
Before the new one: say the Russian for nausea, then the Russian for dizziness.
\end{quote}

\subsection*{You'll want to know: \ru{обморок}}
\textbf{\ru{обморок}} (\emph{óbmorok}) — \textquotedblleft{}a faint\textquotedblright{}.

\begin{quote}
\ru{В} \ru{душном} \ru{зале} \ru{одна} \ru{женщина} \ru{упала} \ru{в} \ru{обморок}.
\end{quote}

\begin{quote}
\emph{V dúshnom zále odná zhénshchina upála v óbmorok.}
\end{quote}

\begin{quote}
In the stuffy hall a woman fainted.
\end{quote}

Masculine. \emph{\ru{Упасть} \ru{в} \ru{обморок}} — to faint.

\begin{cousinweb}
\emph{\ru{обморок}} is built on \emph{\ru{об}-} and \emph{\ru{морок}}, which EtymDB glosses 'darkness, cloud, fog': fainting is a darkening of the eyes.
\end{cousinweb}

\begin{grammarlens}[title={What you've built}]
One more word for telling, explaining and describing at B1.
\end{grammarlens}

\subsection*{Your turn}
\begin{itemize}
\raggedright
  \item \emph{Say it:} \emph{óbmorok}
  \item \emph{Say it:} \emph{óbmorok}, once more
  \item \emph{Say it:} a sentence of your own with \emph{óbmorok}
  \item \emph{Recall:} say the Russian for nausea, then the Russian for dizziness, then say \emph{óbmorok} again
\end{itemize}

\begin{tcolorbox}[breakable,colback=teal!4,colframe=teal!35!black,title={Before you move on}]
What does \ru{обморок} mean? (\textquotedblleft{}A faint\textquotedblright{}.) Say it once more.
\end{tcolorbox}
\section[\texorpdfstring{\ru{знобить} (znobít)}{знобить (znobít)}]{\ru{знобить} (znobít) — to have the shivers (impersonal)}

\label{lesson:RU-C485-znobit}



\begin{quote}
Before the new one: say the Russian for dizziness, then the Russian for a faint.
\end{quote}

\subsection*{You'll want to know: \ru{знобить}}
\textbf{\ru{знобить}} (\emph{znobít}) — \textquotedblleft{}to have the shivers (impersonal)\textquotedblright{}.

\begin{quote}
\ru{У} \ru{меня} \ru{температура}, \ru{и} \ru{меня} \ru{знобит}.
\end{quote}

\begin{quote}
\emph{U menyá temperatúra, i menyá znobít.}
\end{quote}

\begin{quote}
I've got a temperature and I've got the shivers.
\end{quote}

Impersonal: \emph{\ru{меня} \ru{знобит}}. Imperfective only.

\begin{cousinweb}
\emph{\ru{знобить}} shares its root with \emph{\ru{озноб}} 'a chill'; EtymDB has no entry for either, and the origin is not settled.
\end{cousinweb}

\begin{grammarlens}[title={What you've built}]
One more word for telling, explaining and describing at B1.
\end{grammarlens}

\subsection*{Your turn}
\begin{itemize}
\raggedright
  \item \emph{Say it:} \emph{znobít}
  \item \emph{Say it:} \emph{znobít}, once more
  \item \emph{Say it:} a sentence of your own with \emph{znobít}
  \item \emph{Recall:} say the Russian for dizziness, then the Russian for a faint, then say \emph{znobít} again
\end{itemize}

\begin{tcolorbox}[breakable,colback=teal!4,colframe=teal!35!black,title={Before you move on}]
What does \ru{знобить} mean? (\textquotedblleft{}To have the shivers (impersonal)\textquotedblright{}.) Say it once more.
\end{tcolorbox}
